\documentclass{article}
\usepackage{amsmath}
\usepackage{graphicx}
\graphicspath{ {images/} }
\usepackage[spanish]{babel}
\usepackage{bbm}
\usepackage{amsthm}
\usepackage{authblk}
\usepackage{hyperref}
\usepackage[utf8]{inputenc}
\usepackage{mathrsfs}
\usepackage{amsfonts}
\usepackage{amssymb}
\usepackage{enumerate}
\usepackage[left=3cm,right=2cm,top=2cm,bottom=2cm]{geometry}
\usepackage{epsfig}
 \renewcommand{\baselinestretch}{0.93}
 \thispagestyle{empty}
 \pagestyle{empty}
\setlength{\textheight}{53pc} \setlength{\textwidth} {36pc}
\setlength{\topmargin}{-4mm}
\usepackage{multicol}
\newcommand*{\email}[1]{%
    \normalsize\href{mailto:#1}{#1}\par
    }
\setlength{\columnsep}{1cm}
\newcommand{\N}{\mathbb{N}}
\newcommand{\calM}{\cal{M}}
\newcommand{\calP}{\cal{P}}
\newcommand{\F}{\mathbb{F}}
\newcommand{\R}{\mathbb{R}}
\newcommand{\Z}{\mathbb{Z}}
\newcommand{\Q}{\mathbb{Q}}
\newcommand{\C}{\mathbb{C}}
\newcommand{\bbH}{\mathbb{H}}


\theoremstyle{definition}
\newtheorem{ejer}{Ejercicio}
\author{Marcos Morales}
\affil{Universidad Finis Terrae}
\title{Primera Semana\\}



	


\begin{document}


\begin{picture}(400, 75)
   \put(-15,15){\includegraphics[width=5cm]{UFT.png}}
\end{picture}



\begin{center}
\textbf{\Large{Clase Práctica}} \\
\end{center}

\setlength{\parindent}{0cm}



\begin{center}
\textbf{Cálculo Vectorial}
\end{center}

\hspace{100cm} \\ 


\textbf{Ejercicio 1:} Calcule las siguientes integrales dobles \\

\begin{enumerate}
    \item $\displaystyle \iint_{[0,2]\times[0,1]} (1+2x+2y)dA$ \\
    \item $\displaystyle \iint_{[0,2]\times[x,2]} x\sin(xy)dA$ \\
    \item $\displaystyle \int_1^2\int_x^{x\sqrt{3}} xy dydx$ \\
    \item $\displaystyle \int_0^1 \int_{\sqrt{1-y^2}}^{-\sqrt{1-y^2}}dxdy$ \\
\end{enumerate}


\textbf{Ejercicio 2:} Determine el volumen bajo la superficie $f(x,y) = x^2e^{xy}$, sobre el triángulo de vértices $(0,0)$, $(0,1)$ y $(2,1)$. \\

\textbf{Ejercicio 3:} Determine el volumen bajo la superficie  $f(x,y)=1$ sobre la región limitada por el eje $Y$, el eje $X$, la recta $x=1/2$ y la recta $y=1-x$ \\

\textbf{Ejercicio 4:} Calcule 


\begin{align*}
    \iint_R e^{\frac{y-x}{x+y}}dA.
\end{align*}

Donde $R$ es la región del plano acotada por el eje $X$, el eje $Y$ y la recta $x+y+1=0$.

\end{document} 